\lysh{38917}{4}

\begin{alist}
\item
\begin{rlist}
\item Η δέσμη νερού φτάνει στο έδαφος όταν το ύψος $\upsilon(x) = -0,13x^2 + x + 3$ μηδενίζεται. Έτσι διαδοχικά έχουμε:
$\upsilon(x) = 0 \Leftrightarrow -0,13x^2 + x + 3 = 0 \Leftrightarrow 13x^2 - 100x - 300 = 0$. Η εξίσωση αυτή είναι $2$ου βαθμού με διακρίνουσα $\Delta = (-100)^2 - 4 \cdot 13 \cdot (-300) = 10000 + 15600 = 25600$ και ρίζες
$x_{1,2} = \dfrac{-(-100) \pm \sqrt{25600}}{2 \cdot 13} = \dfrac{100 \pm 160}{26}$, από όπου προκύπτει $x = 10$ (δεκτή), ή $x = -\dfrac{30}{13}$ (που απορρίπτεται επειδή $x > 0$).
Επομένως το νερό φτάνει στο έδαφος σε απόσταση $10$ μέτρων από το όχημα.

\item Τη στιγμή που η δέσμη βρίσκεται σε ύψος $4$ μέτρων ισχύει:
$\upsilon(x) = 4 \Leftrightarrow -0,13x^2 + x + 3 = 4 \Leftrightarrow -0,13x^2 + x - 1 = 0 \Leftrightarrow 13x^2 - 100x + 100 = 0$. Η τελευταία εξίσωση είναι $2$ου βαθμού με διακρίνουσα $\Delta = (-100)^2 - 4 \cdot 13 \cdot 100 = 10000 - 5200 = 4800$ και ρίζες
$x_{1,2} = \dfrac{-(-100) \pm \sqrt{4800}}{2 \cdot 13} \simeq \dfrac{100 \pm 69}{26}$, από όπου προκύπτει $x_1 = \dfrac{31}{26} \simeq 1,2$ και $x_2 = 6,5$.

Οπότε για να φθάσει στο έδαφος η δέσμη νερού, δηλαδή στα $10\text{ m}$ από το πυροσβεστικό, θα πρέπει να διανύσει άλλα $8,8\text{ m}$, αν βρίσκεται στη θέση $x_1$ και άλλα $3,5\text{ m}$, αν βρίσκεται στη θέση $x_2$.
\end{rlist}

\item
\begin{rlist}
\item Επειδή η φωτιά βρίσκεται σε ύψος $9$ μέτρων, για να τη φτάσει η δέσμη νερού που θα εκτοξεύσει το πυροσβεστικό, θα πρέπει το πυροσβεστικό να τοποθετηθεί στις θέσεις $x$, για τις οποίες ισχύει $h(x) \ge 9$ από όπου ισοδύναμα έχουμε:
$-x^2 + 10x \ge 9 \Leftrightarrow x^2 - 10x + 9 \le 0$, που αποτελεί ανίσωση $2$ου βαθμού με διακρίνουσα $ \Delta = (-10)^2 - 4 \cdot 1 \cdot 9 = 100 - 36 = 64$ και ρίζες
$x_{1,2} = \dfrac{-(-10) \pm \sqrt{64}}{2 \cdot 1} = \dfrac{10 \pm 8}{2}$ ή $x_1 = 1, x_2 = 9$. Το τριώνυμο $x^2 - 10x + 9$ είναι ετερόσημο του $\alpha = 1$, δηλαδή αρνητικό ή ίσο με το μηδέν στο κλειστό διάστημα μεταξύ των ριζών του. Επομένως, ισχύει $h(x) \ge 9$ για $x \in [1, 9]$, δηλαδή το πυροσβεστικό θα πρέπει να απέχει οριζόντια απόσταση $1\text{ m}$ έως $9\text{ m}$ από τη φωτιά, για να τη φτάνει.

\item Αρκεί να αποδείξουμε ότι για οποιαδήποτε οριζόντια απόσταση $x$ της δέσμης από το πυροσβεστικό όχημα ισχύει ότι:
$h(x) \le 25$ ή $-x^2 + 10x \le 25 \Leftrightarrow x^2 - 10x + 25 \ge 0 \Leftrightarrow (x - 5)^2 \ge 0$, που είναι προφανές για οποιαδήποτε τιμή του $x$. Συνεπώς, οποιαδήποτε φωτιά βρίσκεται σε ύψος μεγαλύτερο των $25$ μέτρων δεν μπορεί να τη φτάσει το πυροσβεστικό όχημα με τη συγκεκριμένη ρύθμιση της δέσμης νερού.
\end{rlist}
\end{alist}

